\begin{abstract}
%v2, some more technical info from Roberto, keeping Ray's trimmings.
Effective robotic systems for long-horizon human-robot collaboration must adapt to a wide range of human partners, whose physical behavior, willingness to assist, and understanding of the robot's capabilities may change over time. This demands a tightly coupled communication loop that grants both agents the flexibility to propose, accept, or decline requests as they coordinate toward completing the task effectively. 
% We apply a \textbf{M}ixed-\textbf{I}nitiative dialog paradigm to \textbf{Co}llaborative human-ro\textbf{Bot} teaming and propose \textbf{\methodname{}}, a system that handles the common scenario where both agents, using natural language, take initiative in formulating, accepting, or rejecting proposals on who can best complete different steps of a task.
We propose \textbf{\methodname{}}, a system that enables the human and robot, both using natural language, to take initiative in formulating, accepting, or rejecting proposals on who can best complete different steps of a task.
To handle diverse, task-directed dialog, and find successful collaborative strategies that minimize human effort, MICoBot makes decisions at three levels: (1) a meta-planner considers human dialog to formulate and code a high-level collaboration strategy, (2) a planner optimally allocates the remaining steps to either agent based on the robot's capabilities (measured by a simulation-pretrained affordance model) and the estimated human's willingness to help, and (3) an action executor decides the low-level actions to perform or words to say to the human.
% Our extensive evaluations in simulation and real-world---on a physical robot with 18 unique human participants over 27 hours---demonstrate our method's effectiveness at collaborating with diverse human users, yielding significantly improved task success and user experience than a pure LLM baseline and other agent allocation models.
In physical robot trials with 18 unique human participants, MICoBot significantly improves task success and user experience over a pure LLM baseline and standard agent allocation models.
% We demonstrate MICoBot's effectiveness at collaborating with diverse human users through extensive evaluations in simulation and the real-world---on a physical robot with 18 unique human participants---yielding significantly improved task success and user experience than a pure LLM baseline and other agent allocation models.
See additional videos and materials at our project site.
% \footnote{https://mico-bot.github.io\label{proj-site}}
\footnote{https://robin-lab.cs.utexas.edu/MicoBot/\label{proj-site}}.
% \url{https://robin-lab.cs.utexas.edu/MicoBot/}.
% \url{https://mico-bot.github.io}.


% % v1, no technical info. trimmed by Ray.
% Effective robotic systems that work with humans on long-horizon tasks must adapt to a wide distribution of human collaborators with variable physical behavior, willingness to help, and understanding of the robot's abilities. This requires a tight human-robot communication loop that affords each agent latitude to propose, accept, and reject requests  when deciding how to best complete a task together.
% %Prevailing paradigms in human-chatbot and human-robot systems involve a single agent driving the topic and content of the dialog exchanges while the other fulfills all the requests of the former. These settings constrict the communication flow between the two agents and can greatly restrict the seamlessness and adaptability of human-robot teamwork.
% %Recently, mixed-initiative dialog systems have been proposed to 
% We apply a mixed-initiative dialog paradigm to human-robot collaboration, where both agents, using natural language, take initiative in formulating, accepting, or rejecting proposals on who can best complete different steps of a task.
% Our approach leverages the code generation and language understanding abilities of LLMs while compensating for their lack of knowledge of the robot's capabilities through a learned affordance-based approach.
% % The resulting system computationally optimizes how to split up a task depending on the robot's capabilities and the human's inferred future helpfulness.
% Our evaluations in simulation and the real-world
% %---on a physical robot with 18? unique human participants---
% demonstrate the ability of our method to effectively collaborate despite diverse human behavior, yielding greater task success than a pure LLM baseline and simpler agent allocation models.
\end{abstract}